\documentclass[conference]{IEEEtran}
\usepackage[T1]{fontenc}
\usepackage{amsmath,amssymb,amsthm,mathtools}
\usepackage{booktabs,microtype,array}
\usepackage{graphicx,pgfplots}
\usepackage[hidelinks]{hyperref}
\pgfplotsset{compat=1.18}
\hypersetup{pdftitle={Certified Accounting for Adaptive Shuffled Gaussian Training},pdfauthor={Samuel Mausberg}}
\newtheorem{theorem}{Theorem}
\newtheorem{lemma}[theorem]{Lemma}
\newtheorem{proposition}[theorem]{Proposition}
\newtheorem{corollary}[theorem]{Corollary}
\newcommand{\E}{\mathbb E}
\newcommand{\R}{\mathbb R}
\newcommand{\N}{\mathcal N}
\newcommand{\Prb}{\mathbb P}
\newcommand{\HS}[3]{H_{#1}(#2,#3)}
\newcommand{\pos}[1]{\left(#1\right)_+}
\newcommand{\TV}{\operatorname{TV}}
\newcommand{\Bin}{\operatorname{Bin}}
\newcommand{\eps}{\varepsilon}
\newcommand{\pair}{\mathrm{pair}}
\newcommand{\dd}{\,\mathrm d}
\title{Certified Accounting for Adaptive\\Shuffled Gaussian Training}
\author{\IEEEauthorblockN{Samuel Mausberg}\IEEEauthorblockA{Independent Researcher}}
\begin{document}
\pagestyle{plain}
\bstctlcite{BSTcontrol}
\maketitle
\thispagestyle{plain}
\begin{abstract}
Many differentially private training pipelines shuffle the data once per epoch and read fixed-size batches, but account for privacy as if the batches were Poisson subsampled. This mismatch can understate the cost. We give a certified accountant for multi-epoch shuffled training with Gaussian noise. It remains valid when each query may depend on all earlier releases and the items are heterogeneous. A null-clone reduction hides the differing item among a random number of candidates. We bound the likelihood sum at each candidate count and merge the resulting bounds into one finite rational dominating pair that composes across epochs. With 1000 batches per epoch, five epochs, noise multiplier $1$ and $\delta=10^{-8}$, the worst-case epsilon of the full adaptive transcript lies in $(6.423,7.97]$, while matched Poisson sampling is certified at $0.507$. We also show that the pair conjectured by Chua et al.\ to dominate all adaptive queries does not. The dominating-pair construction and the counterexample are verified in Lean~4, along with supporting lemmas for the reduction and the certificate.
\end{abstract}

\section{Introduction}
A training pipeline that shuffles its data and then reads consecutive batches uses each item exactly once per epoch. This rule is attractive to implement. It is also part of the privacy mechanism: a batch containing the differing item contains one fewer common item, and that missing contribution can help identify its position. Chua et al.~\cite{chua2024,scalable} showed that replacing this sampling rule by Poisson subsampling in the analysis can materially understate the privacy cost of the adaptive-query transcript.

With $T$ batches per epoch and noise multiplier $\sigma$, the candidate worst case of Chua et al.\ is
\begin{equation}\label{eq:chua}
 P_S=\frac1T\sum_{j=1}^T\N(2e_j,\sigma^2 I_T),\qquad
 Q_S=\frac1T\sum_{j=1}^T\N(e_j,\sigma^2 I_T).
\end{equation}
Here $e_j$ is the $j$th standard basis vector of $\R^T$. A scalar fixed query attains this pair: the differing item has value $+1$, every common item has value $-1$, and the neighboring input replaces the differing item by zero. Conjecture 3.2 of~\cite{chua2024} proposes that this pair tightly dominates the shuffled mechanism for all adaptive query methods. It does not: at every real $\eps$ an adaptive final query, chosen for that $\eps$, strictly increases the forward divergence above the pair's value, and at every $\eps\ge0$ the same query also increases the reverse divergence. Appendix~\ref{app:reflection} gives the proof. The pair nevertheless remains a useful attainable lower benchmark.

We give upper and lower certificates at practical settings. For $T=1000$, five epochs, and $\delta=10^{-8}$, the unrestricted adaptive epsilon lies in
\begin{equation}\label{eq:brackets}
 (6.423,7.97]\quad(\sigma=1),\qquad
 (0.338,0.71]\quad(\sigma=2).
\end{equation}
The lower endpoints follow from an explicit statistic of the independently repeated pair~\eqref{eq:chua}. The upper endpoints apply to heterogeneous items, every batch size, and arbitrary adaptation within and between epochs. Revealing the permutations gives the deterministic-batching upper endpoints $14.55021$ and $6.54266$.

The upper endpoints come from a sharper evaluation of a clone reduction. After the common kernels acquire an approximate null component, the distinguished item is hidden among $1+\Bin(T-1,p)$ candidates. High moments of the likelihood ratio obscure the effect of small candidate counts. We bound the conditional likelihood sum at each relevant count, average those bounds over the count, and only then build a composable experiment. At five epochs the upper bound falls from $9.21$, under a moment/Laplace evaluation of the same reduction, to $7.97$ at $\sigma=1$, and from $0.83$ to $0.71$ at $\sigma=2$.

A matched Poisson calculation puts the brackets on a familiar scale. With independent inclusion probability $1/T$, the same noise and $T$ updates per epoch, Poisson accounting certifies $0.507$ and $0.182$ at five epochs. The shuffled lower endpoints exceed those values by factors greater than $12.6$ and $1.85$.

Section~\ref{sec:model} states the model and the main guarantee. Section~\ref{sec:clone} bounds each conditional epoch through null clones and explicit candidate counts, Section~\ref{sec:composition} turns those bounds into one composable pair, and Section~\ref{sec:evaluation} reports the certified comparisons. Section~\ref{sec:lean} describes the Lean~4 formalization and lists the steps that rest on the written proofs.

\section{Related work and contribution}
Cloning is the established amplification principle behind the reduction. Feldman, McMillan, and Talwar~\cite{fmt} analyze sequentially adaptive randomizers using binomially distributed clones. Their approximate-DP analysis also replaces kernels at a total-variation cost; Lemma C.1 of the third arXiv version records the $(1+e^\eps)$ transfer charge. We specialize these ingredients to a fixed null Gaussian and keep the distinguished Gaussian kernel intact.

Composition through hypothesis-testing tradeoffs and dominating pairs is standard~\cite{gdp,zhu}. Gopi, Lee, and Wutschitz~\cite{glw} develop numerical privacy-loss composition with controlled approximation error, and Doroshenko et al.~\cite{dots} construct valid discrete experiments from privacy-profile information. Our rational testing envelope is a dual form of this discretization principle. The numerical contribution is a directed candidate-count calculation and its use in a certified shuffled-training accountant.

Feldman and Shenfeld~\cite{fs2025,fs2026} give reductions and efficient composable accounting for independent random allocation. Their likelihood-sum computations are closely related to the conditional experiment used here. Fixed-capacity shuffling requires an additional reduction before those experiments apply: conditioning on all common assignments identifies the null slot. Section~\ref{sec:capacity} tracks this distinction through the proof. Recent analyses of a fixed randomize-then-shuffle channel~\cite{takagi} and of an all-noise versus hidden-shift experiment for random allocation~\cite{vandijk} also depend on the observation model in this way.

For balls-and-bins sampling, which is random allocation with one step per item, Chua et~al.~\cite{ballsbins} prove that the all-noise versus hidden-shift pair with $T$ positions is tightly dominating for all adaptive query methods, and they estimate its divergences by Monte Carlo sampling with high-probability guarantees. After a total-variation charge for modifying the common kernels, Lemma~\ref{lem:clone} reduces each conditional epoch to the same pair with the random number $K$ of positions, and we bound the resulting divergences with directed arithmetic.

The dynamic-shuffling lower construction in \emph{Scalable DP-SGD}~\cite[Sec.~3, before Thm.~3.7]{scalable} attains $(P_S^{\otimes E},Q_S^{\otimes E})$. We use this construction and compute certified lower bounds by coarsening each epoch's coordinate maximum. The persistent-shuffling bound in the same paper, printed as Theorem 3.5 in both the arXiv (v1) and NeurIPS versions, omits $e^\eps$ before $Q(\Gamma_C)$, while the paragraph before it gives the correct hockey-stick expression. Our lower bounds use the hockey-stick expression, and we make no assertion about the implementation used in that work.

Two further works measure the cost of accounting for the wrong sampler. Lebeda et~al.~\cite{lebeda} show that fixed-size batches drawn independently without replacement need twice the noise multiplier of Poisson subsampling under add or remove adjacency, so that a Poisson accountant reporting $\eps\approx1$ can correspond to $\eps>10$. They also show that self-composing a worst-case dataset pair of one subsampled step can understate the privacy loss of the composition; Theorem~\ref{thm:main} avoids that step by composing a pair that dominates every conditional epoch. Annamalai et~al.~\cite{annamalai} audit DP-SGD with shuffling and find empirical leakage up to four times the Poisson guarantees reported for state-of-the-art models. Their audits give statistical lower bounds under particular threat models, while the brackets in~\eqref{eq:brackets} are certified on both sides.

\section{Model and main guarantee}\label{sec:model}
Let $b$, $T$, $E$ and $d$ be positive integers and $\sigma>0$. Put $n=bT$ and fix items $x_1,\ldots,x_n$. Training runs for $E$ epochs. Every epoch draws a fresh independent uniform permutation and divides it into $T$ consecutive batches of size $b$; $B_t$ is the batch read at step $t\le ET$. From all preceding releases, a measurable analyst chooses $\psi_t$ with values in $\R^d$ satisfying
\[
 \|\psi_t(x)\|_2\le1,\qquad \psi_t(\bot)=0.
\]
The release is
\begin{equation}\label{eq:release}
 Y_t=\sum_{i\in B_t}\psi_t(x_i)+\sigma Z_t,
 \qquad Z_t\stackrel{\mathrm{iid}}\sim\N(0,I_d).
\end{equation}
Neither batch identities nor permutation state is exposed. Analyst randomness is independent of the private input and can be conditioned on. Neighbors replace one item by $\bot$ while retaining its slot. The common items may be arbitrary.

Write $P$ for the real-item law and $Q$ for its null neighbor. For real $\eps$, define
\[
 H_\eps(P,Q)=\sup_A\{P(A)-e^\eps Q(A)\}.
\]
A pair $(R,S)$ dominates $(P,Q)$ if its forward divergence is no smaller at every real threshold. Because $H_\eps(Q,P)=1-e^\eps+e^\eps H_{-\eps}(P,Q)$ for probability measures, domination also controls the reverse divergence. Let $\eps^*_{E,T,\sigma}(\delta)$ be the smallest $\eps$ such that $H_\eps(P,Q)\le\delta$ and $H_\eps(Q,P)\le\delta$ for the $E$-epoch transcript laws of every admissible analyst, dimension, batch size, dataset and zero-out neighbor. Let $\eps_{\pair,E}(\delta)$ be the smallest $\eps$ with $\max\{H_\eps(P_S^{\otimes E},Q_S^{\otimes E}),H_\eps(Q_S^{\otimes E},P_S^{\otimes E})\}\le\delta$ for the pair~\eqref{eq:chua}.

An epoch accountant must remain valid after conditioning on the earlier training history. Once that property is established, a single finite experiment can be reused through the run. The following theorem states the resulting guarantee; the next sections construct and verify its inputs.

\begin{theorem}[Certified adaptive accounting]\label{thm:main}
Suppose each conditional epoch pair obeys finitely many inequalities
\[
 H_{\log k_j}(P,Q)\le d_j^+,\qquad
 H_{\log k_j}(Q,P)\le d_j^-,\qquad k_j\ge1,
\]
with rational $k_j$ and $d_j^\pm$, and $0\le d_j^\pm\le1$. For $a\in[0,1]$, set
\begin{equation}\label{eq:envelope}
 f(a)=\max\left\{0,\ \max_j(1-d_j^+-k_ja),\
                    \max_j\frac{1-d_j^- -a}{k_j}\right\}.
\end{equation}
There is a finite rational pair $(R,S)$ with testing tradeoff $f$ that dominates every conditional epoch. Its $E$-fold product dominates the entire adaptive transcript for $E$ fresh reshuffles.

The certified inputs supplied with this paper give, at $T=1000$ and $\delta=10^{-8}$,
\begin{align*}
 6.423&<\eps_{\pair,5}(\delta)\le\eps^*_{5,1000,1}(\delta)\le7.97,\\
 0.338&<\eps_{\pair,5}(\delta)\le\eps^*_{5,1000,2}(\delta)\le0.71.
\end{align*}
In each line the product pair uses the same $\sigma$ as the mechanism.
\end{theorem}

Scaling every clipped gradient and the Gaussian noise by a public $C>0$ leaves the guarantee unchanged, so $\sigma$ is the usual noise multiplier, the noise standard deviation divided by $C$. Dividing the sum by the fixed or expected batch size is postprocessing. Thus the upper accountant applies to the usual clipped-gradient implementation of DP-SGD under the specified sampling rule.

\section{Null clones and the candidate count}\label{sec:clone}
A common release that has the null law cannot identify which candidate slot holds the differing item. We create such releases in an auxiliary experiment and charge the change before composing epochs. The construction is conditional on an arbitrary incoming history.

\subsection{One unrevealed item per batch}\label{sec:singleton}
Let $i_\star$ be the differing index. Choose $T-1$ further indices uniformly from the other $n-1$; together with $i_\star$ they are the representatives. Independently allocate the remaining $n-T$ indices uniformly into labeled filler groups $G_1,\ldots,G_T$ of size $b-1$, and independently draw a uniform bijection $\pi$ from batches to representatives, so that batch $t$ is $G_t\cup\{\pi(t)\}$. A triple produces a fixed ordered partition exactly when $i_\star$ represents its own block and every other block supplies one of its $b$ elements as representative. That choice determines the filler groups and $\pi$, so each ordered partition arises from $b^{T-1}$ equally likely triples, and the batches form a uniform ordered partition.

Reveal $W$, consisting of the representative set and the filler groups, and keep $\pi$ hidden. The revealed $W$ does not depend on the data, and every filler is a common item. Given $W$ and the incoming history, the filler sum $c_t=\sum_{i\in G_t}\psi_t(x_i)$ is a measurable function of $Y_{<t}$. The map $Y_t\mapsto Y_t-c_t(Y_{<t})$ has the recursive measurable inverse $Y_t=Y'_t+c_t(Y_{<t})$, so it changes neither ordered divergence, and it leaves an adaptive singleton shuffle of the $T$ representatives with the same $\sigma$. Because $W$ has the same law under both hypotheses, $H_\eps$ of the joint law of $(W,Y)$ is the $W$-average of the conditional divergences in either order, and data processing bounds the divergence of $Y$ alone by it.

\subsection{The conditional null-clone inequality}
The kernel modification has a scalar cost even when the query dimension is large. Fix $p\in(0,1)$, define $h=\log(1/p)$ and
\begin{align}
 \gamma_\sigma(p)&=p\Phi\left(-\sigma h+\frac1{2\sigma}\right)
 -\Phi\left(-\sigma h-\frac1{2\sigma}\right),\label{eq:gamma}\\
 \beta_p&=(T-1)\gamma_\sigma(p).\nonumber
\end{align}
This cost applies uniformly to every common query vector in the unit ball, which permits the same argument along adaptive histories.

\begin{lemma}[Null-clone reduction]\label{lem:clone}
Fix $p\in(0,1)$ and let $K=1+\Bin(T-1,p)$, independently of iid standard normals $Z_i$, and put
\[
 X_i=e^{Z_i/\sigma-1/(2\sigma^2)},\qquad
 L_m=\frac1m\sum_{i=1}^mX_i.
\]
For every conditional neighboring epoch pair and $k\ge1$,
\begin{align}
 H_{\log k}(P,Q)&\le\E(L_K-k)_++(1+k)\beta_p,\label{eq:clonef}\\
 H_{\log k}(Q,P)&\le\E(1-kL_K)_++(1+k)\beta_p.\label{eq:cloner}
\end{align}
\end{lemma}
\begin{proof}
Use the singleton reduction above, and let $g_u$ be the density of $\N(u,\sigma^2I_d)$. Set
\[
 \eta_u=\int(pg_0-g_u)_+,\qquad
 r_u=\frac{(g_u-pg_0)_+}{1-p+\eta_u}.
\]
The modified density $\widetilde g_u=pg_0+(1-p)r_u$ integrates to one and satisfies $\TV(g_u,\widetilde g_u)=\eta_u$. A Gaussian displacement of norm at most one is a postprocessing of a scalar unit displacement, so $\eta_u\le\gamma_\sigma(p)$; at a unit scalar displacement, the Gaussian likelihood-ratio threshold evaluates the integral as~\eqref{eq:gamma}.

Modify only the common kernels. Conditionally on $\pi$, exactly $T-1$ history-dependent kernels change, each by at most $\gamma_\sigma(p)$ in total variation at every history, because equal prefixes induce the same query vector. A sequential coupling up to its first disagreement, averaged over $\pi$, gives
\[
 \TV(P,\widetilde P),\ \TV(Q,\widetilde Q)\le\beta_p.
\] Transferring a hockey-stick inequality through the two errors costs at most $(1+k)\beta_p$.

To realize the modified mechanism, attach an independent Bernoulli-$p$ coin to each common index. A successful coin emits $g_0$; a failure emits $r_u$. Reveal the clone indices and the positions of all nonclones. The differing index is uniform among the remaining $K$ positions. Every other candidate emits the null law.

The conditional experiment is a causal postprocessing of
\[
 P_K=\frac1K\sum_{j=1}^K\N(e_j,\sigma^2I_K),\qquad
 Q_K=\N(0,\sigma^2I_K).
\]
At a candidate position with current differing-item query vector $a$, transform its scalar input $V_j$ to
\[
 aV_j+\sigma(I_d-aa^\top)^{1/2}Z'_j.
\]
Here $a=\psi_t(x_{i_\star})$ is computed from the simulated prefix and $Z'_j\sim\N(0,I_d)$ is fresh. Since $\|a\|_2\le1$, the matrix $I_d-aa^\top$ is positive semidefinite and its square root is continuous in $a$, so each step is a measurable kernel. A scalar input with mean $0$ gives $g_0$ and one with mean $1$ gives $g_a$, both with covariance $\sigma^2I_d$. The simulator uses $x_{i_\star}$ but not the hypothesis, so the same kernel acts on both laws. Nonclones are simulated from their residual kernels. Conditional on the distinguished position, future input noises remain independent of the simulated prefix, so the simulation permits feedback.

Under $Q_K$ the likelihood ratio is $L_K$. Data processing, followed by averaging the common revelation and paying the TV transfer, proves both inequalities.
\end{proof}

\subsection{Keeping small candidate counts explicit}
A moment bound averages different candidate counts before resolving the likelihood tail. That can be expensive when a small count and a large likelihood ratio occur together. We retain the count and compute the two conditional positive-part expectations directly.

Write $S_m=\sum_{i=1}^mX_i$ and define
\[
 F_m(t)=\E(S_m-t)_+,\qquad J_m(t)=\E(t-S_m)_+,
\]
a call and a put on $S_m$.
They obey the scalar recurrences
\begin{equation}\label{eq:recurrences}
 F_m(t)=\E F_{m-1}(t-X),\qquad
 J_m(t)=\E J_{m-1}(t-X),
\end{equation}
with $F_0(t)=(-t)_+$ and $J_0(t)=t_+$. Both are convex in $t$. Convexity permits a finite upper computation using endpoint spreads in $X$ and chords between threshold nodes.

Let $g_\sigma(\eps)$ be the sensitivity-one Gaussian profile,
\begin{equation}\label{eq:gaussian}
 g_\sigma(\eps)=\Phi\left(-\sigma\eps+\frac1{2\sigma}\right)
 -e^\eps\Phi\left(-\sigma\eps-\frac1{2\sigma}\right).
\end{equation}
This profile bounds every conditional epoch in both orders. Once the permutation is revealed, the differing item enters one round, where given the prefix the two laws are $\N(c,\sigma^2I_d)$ and $\N(c+a,\sigma^2I_d)$ with $\|a\|_2\le1$; every other round has the same kernel under both hypotheses, so both ordered divergences at $\log k$ are at most $g_\sigma(\log k)$. Conditionally on $K=m$, Jensen's inequality gives $\E(L_m-k)_+\le\E(X-k)_+=g_\sigma(\log k)$ and $\E(1-kL_m)_+\le g_\sigma(\log k)$, which covers candidate counts beyond the finite convolution horizon.

\begin{proposition}[Finite candidate-count bound]\label{prop:counts}
Suppose $1\le M\le T$ and certified functions satisfy $\overline F_m\ge F_m$, $\overline J_m\ge J_m$ for $m\le M$. Put $w_m=\Prb(K=m)$, $g=g_\sigma(\log k)$, and
\[
\begin{aligned}
 c_m^+&=\min\{g,\overline F_m(mk)/m\},\\
 c_m^-&=\min\{g,k\overline J_m(m/k)/m\}.
\end{aligned}
\]
The respective conditional epoch divergences are at most
\begin{equation}\label{eq:finite-count}
 \sum_{m=1}^Mw_mc_m^\pm+
 \left(1-\sum_{m=1}^Mw_m\right)g+(1+k)\beta_p.
\end{equation}
Taking the minimum with $1$ and $g$ remains valid. Different thresholds and orientations may use different $p$.
\end{proposition}
\begin{proof}
Conditioning on $K=m$ gives $F_m(mk)/m$ in the forward direction and $kJ_m(m/k)/m$ in reverse. Each is at most $g$ by Jensen's inequality. Substitute the finite bounds into Lemma~\ref{lem:clone}; use $g$ for the residual candidate probability. The choice of $p$ changes only the comparison experiment.
\end{proof}

\subsection{Why fixed capacity matters}\label{sec:capacity}
The independent-allocation reduction of Feldman and Shenfeld~\cite{fs2025} (Theorem 3.1 in the NeurIPS version and arXiv v3--v4, Lemma 3.1 in v1--v2) conditions on the common items' assignments without revealing the differing item's assignment. Under fixed-capacity shuffling, the same revelation identifies the batch with only $b-1$ common items and therefore the differing slot.

An inert item can be equivalent to removal under independent allocation. Under shuffling it still occupies capacity. For the antipodal query, the two means are $-b\mathbf1+2e_J$ and $-b\mathbf1+e_J$, yielding~\eqref{eq:chua} after translation. The all-noise null experiment appears in Lemma~\ref{lem:clone} only after paying for the modification and replacing $T$ candidate positions by the random count $K$. Appendix~J of the third arXiv version of~\cite{vandijk} makes the same observation for its scalar all-noise pair, which does not dominate~\eqref{eq:chua}.

\section{An experiment that composes}\label{sec:composition}
A forward privacy inequality limits the probability of accepting the real-item hypothesis. A reverse inequality supplies another limit by applying it to the complementary decision. Their intersection is the testing envelope in Theorem~\ref{thm:main}, so threshold-specific clone choices can be used together.

\begin{proof}[Proof of the composition claim in Theorem~\ref{thm:main}]
For a randomized test $\varphi$, let
\[
 T_{P,Q}(a)=\inf_{0\le\varphi\le1:\,\E_Q\varphi\le a}
 (1-\E_P\varphi).
\]
Integrating hockey-stick inequalities over test level sets extends them to randomized tests. The forward inequality gives $1-\E_P\varphi\ge1-d_j^+-k_ja$. Applying the reverse inequality to $1-\varphi$ gives $(1-d_j^- -a)/k_j$. Thus $T_{P,Q}\ge f$.

The envelope $f$ is convex and nonincreasing, lies between zero and $1-a$, and has $f(1)=0$. On an interval of length $\Delta a_\ell$ with slope $-t_\ell$, assign an atom the masses
\begin{equation}\label{eq:atoms}
 S_\ell=\Delta a_\ell,\qquad R_\ell=t_\ell\Delta a_\ell.
\end{equation}
Add an $R$-only atom of mass $1-f(0)$. A zero-slope interval is an $S$-only atom. The masses are nonnegative and sum to one under each law.

The likelihood ratios $t_\ell$ do not increase along the envelope. A Neyman--Pearson test first selects the $R$-only atom and then successive intervals, randomizing within an atom as necessary. Its success probability at $S$-mass $a$ is exactly $1-f(a)$. The constructed pair has tradeoff $f$ and dominates every conditional epoch by binary experiment comparison~\cite{gdp,zhu,dots}. Rational input lines have rational intersections, giving rational atom masses.

Each fresh permutation is independent of the incoming transcript. The conditional dominance therefore holds at every history, and the adaptive composition theorem for dominating pairs~\cite[Thm.~10]{zhu} (restated in~\cite[Thm.~2.4]{dots}) gives domination by $(R^{\otimes E},S^{\otimes E})$. The same proof allows different certified pairs across epochs.
\end{proof}

The construction is exact relative to the selected testing inequalities. In the supplied instances every finite atom has privacy loss $j\log(51/50)$ for an integer $j$, because the envelope slopes are the thresholds $k_j$ and their reciprocals. We compose the real-law losses and the reverse-oriented null-law losses separately and charge the singular atoms at infinite loss in each direction.

\section{Certified practical comparisons}\label{sec:evaluation}
We fix $T=1000$, $\delta=10^{-8}$ and noise multipliers $1$ and $2$. Every displayed upper endpoint is verified in both directions. The lower curve certifies the product pair through a fixed coarsening described below. Figures~\ref{fig:sigma1} and~\ref{fig:sigma2} plot the rows of Table~\ref{tab:all}; their lines only join certified markers and carry no guarantee between them.

\begin{table*}[t]
\centering
\caption{Certified epsilon bounds at $T=1000$, $\delta=10^{-8}$. The columns marked $L$ are strict lower bounds from an attainable statistic of the continuous product pair. The columns marked $U$ are certified upper bounds; Deterministic is the Gaussian bound with the permutations revealed. Poisson uses $q=1/1000$ and $1000E$ updates.}
\label{tab:all}
\small
\begin{tabular}{r r r r r r r r r}
\toprule
&\multicolumn{4}{c}{$\sigma=1$}&\multicolumn{4}{c}{$\sigma=2$}\\
\cmidrule(lr){2-5}\cmidrule(lr){6-9}
Epochs&Pair $L$&Shuffle $U$&Poisson $U$&Deterministic $U$&Pair $L$&Shuffle $U$&Poisson $U$&Deterministic $U$\\
\midrule
1&5.106&5.64&0.306&5.77610&0.257&0.47&0.081&2.70761\\
2&5.538&6.49&0.363&8.53379&0.289&0.54&0.114&3.94247\\
3&5.885&7.11&0.414&10.77586&0.310&0.60&0.140&4.92507\\
5&6.423&7.97&0.507&14.55021&0.338&0.71&0.182&6.54266\\
10&7.354&9.46&0.697&22.17344&0.383&0.93&0.260&9.69785\\
20&8.646&11.68&0.982&34.45138&0.438&1.29&0.372&14.55021\\
\bottomrule
\end{tabular}
\end{table*}

\subsection{The shuffled upper accountant}
The likelihood thresholds are $k_j=(51/50)^j$, $0\le j\le555$. We compute $\overline F_m$ and $\overline J_m$ for $m\le M$, with $M=64$ at $\sigma=1$ and $M=160$ at $\sigma=2$. The residual binomial probability receives the Gaussian bound in Proposition~\ref{prop:counts}. At each threshold a numerical search proposes $h=\log(1/p)$ on a rational grid of spacing $1/40$; the selected bound is then re-evaluated with directed arithmetic.

The exact endpoint certificates at the five-epoch upper values are
\begin{equation}\label{eq:numeric-cert}
\begin{array}{c|cc}
(\sigma,\eps)&H_\eps(R^{\otimes5},S^{\otimes5})&H_\eps(S^{\otimes5},R^{\otimes5})\\ \hline
(1,7.97)&\le9856064/10^{15}&\le3805278/10^{15}\\
(2,0.71)&\le8525892/10^{15}&\le5244378/10^{15}.
\end{array}
\end{equation}
Here $(R,S)$ is the certified pair of Theorem~\ref{thm:main} at the given $\sigma$. Every entry is below $10^{-8}$, and the exported JSON files record the exact rationals from which these decimals are rounded.

\subsection{An attainable product-pair lower bound}
Independent reshuffling with the fixed antipodal query attains $(P_S^{\otimes E},Q_S^{\otimes E})$: translation by $b\mathbf1$ is a fixed bijection, so this transcript has exactly the divergences of the product pair, and $\eps_{\pair,E}\le\eps^*_{E,T,\sigma}$. For one epoch let $M_Y=\max_jY_j$ after the translation in~\eqref{eq:chua}. Its CDF under displacement $\theta\in\{1,2\}$ is
\begin{equation}\label{eq:maxcdf}
 \Prb_\theta(M_Y\le c)=
 \Phi((c-\theta)/\sigma)\Phi(c/\sigma)^{T-1}.
\end{equation}
We bin this statistic at thresholds $c=\sigma j/100$, $0\le j\le1200$, including both unbounded end bins. Differences of~\eqref{eq:maxcdf} give the two probabilities of every bin. Their product is a postprocessing of the full continuous pair and therefore provides lower bounds on its divergences.

At five epochs, the forward certificates at the lower endpoints are
\[
 H_{6.423}(P_S^{\otimes5},Q_S^{\otimes5})
 \ge10011846/10^{15},
\]
\[
 H_{0.338}(P_S^{\otimes5},Q_S^{\otimes5})
 \ge10201901/10^{15},
\]
for $\sigma=1$ and $\sigma=2$, respectively. Each exceeds $10^{-8}$. The verifier also evaluates the reverse direction. The binned experiment's own inverse is enclosed in $(6.423,6.426]$ and $(0.338,0.342]$; only their lower endpoints are used for the continuous-pair comparison.

\subsection{Matched Poisson and deterministic accounting}
The matched Poisson mechanism includes every item independently with probability $q=1/T$ at each of $ET$ steps. It has expected batch size $b$ and the same noise multiplier. Its tight stepwise dominating pair~\cite{chua2024},~\cite[Thm.~11]{zhu} is
\begin{equation}\label{eq:poisson}
 A=(1-q)\N(0,\sigma^2)+q\N(1,\sigma^2),\qquad
 B=\N(0,\sigma^2).
\end{equation}
We compute a pessimistic discretization of this pair by spreading each conditional likelihood ratio to its two bin endpoints while preserving its mean under $B$. Convex order proves dominance. Integer convolution then certifies the $ET$-step product, including both directions and all tail deficits.

Deterministic batching composes as a Gaussian shift with displacement $\sqrt E/\sigma$. Directed bisection gives intervals of width $10^{-5}$, whose upper endpoints appear in Table~\ref{tab:all}. At five epochs, the Poisson upper values are $0.507$ and $0.182$, whereas the shuffled lower values are $6.423$ and $0.338$. The ratios $6.423/0.507>12.6$ and $0.338/0.182>1.85$ establish a separation without assuming that either reported upper curve is tight.

\begin{figure}[t]
\centering
\begin{tikzpicture}
\begin{semilogyaxis}[
 width=\columnwidth,height=5.2cm,
 xlabel={Epochs},ylabel={Privacy parameter $\eps$},
 xmin=.7,xmax=20.3,ymin=.2,ymax=45,
 xtick={1,5,10,15,20},ytick={.2,.5,1,2,5,10,20,40},
 yticklabels={0.2,0.5,1,2,5,10,20,40},grid=major,
 tick label style={font=\footnotesize},label style={font=\footnotesize},
 legend columns=2,
 legend style={font=\scriptsize,at={(.5,1.03)},anchor=south,draw=none,fill=none,
 column sep=3pt},
 legend cell align=left]
\addplot+[color=blue!65!black,mark=*,mark size=1.5pt,mark options={draw=blue!65!black,fill=blue!65!black},thick] table[x=epochs,y=new_upper]{comparison_sigma1.dat};\addlegendentry{Adaptive shuffle upper}
\addplot+[color=black,mark=triangle*,mark size=1.8pt,mark options={draw=black,fill=black},densely dashed,thick] table[x=epochs,y=pair_lower]{comparison_sigma1.dat};\addlegendentry{Product-pair lower}
\addplot+[color=black!60,mark=square*,mark size=1.3pt,mark options={draw=black!60,fill=black!60},dashdotted] table[x=epochs,y=deterministic]{comparison_sigma1.dat};\addlegendentry{Deterministic upper}
\addplot+[color=red!65!black,mark=diamond*,mark size=1.8pt,mark options={draw=red!65!black,fill=red!65!black},densely dotted,thick] table[x=epochs,y=poisson]{comparison_sigma1.dat};\addlegendentry{Poisson upper}
\end{semilogyaxis}
\end{tikzpicture}
\caption{Certified bounds for noise multiplier $\sigma=1$, $T=1000$ and $\delta=10^{-8}$, on a logarithmic vertical axis. The product-pair lower markers use the per-epoch maximum statistic.}
\label{fig:sigma1}
\end{figure}

\begin{figure}[t]
\centering
\begin{tikzpicture}
\begin{semilogyaxis}[
 width=\columnwidth,height=5.2cm,
 xlabel={Epochs},ylabel={Privacy parameter $\eps$},
 xmin=.7,xmax=20.3,ymin=.055,ymax=20,
 xtick={1,5,10,15,20},ytick={.1,.2,.5,1,2,5,10,20},
 yticklabels={0.1,0.2,0.5,1,2,5,10,20},grid=major,
 tick label style={font=\footnotesize},label style={font=\footnotesize},
 legend columns=2,
 legend style={font=\scriptsize,at={(.5,1.03)},anchor=south,draw=none,fill=none,
 column sep=3pt},
 legend cell align=left]
\addplot+[color=blue!65!black,mark=*,mark size=1.5pt,mark options={draw=blue!65!black,fill=blue!65!black},thick] table[x=epochs,y=new_upper]{comparison_sigma2.dat};\addlegendentry{Adaptive shuffle upper}
\addplot+[color=black,mark=triangle*,mark size=1.8pt,mark options={draw=black,fill=black},densely dashed,thick] table[x=epochs,y=pair_lower]{comparison_sigma2.dat};\addlegendentry{Product-pair lower}
\addplot+[color=black!60,mark=square*,mark size=1.3pt,mark options={draw=black!60,fill=black!60},dashdotted] table[x=epochs,y=deterministic]{comparison_sigma2.dat};\addlegendentry{Deterministic upper}
\addplot+[color=red!65!black,mark=diamond*,mark size=1.8pt,mark options={draw=red!65!black,fill=red!65!black},densely dotted,thick] table[x=epochs,y=poisson]{comparison_sigma2.dat};\addlegendentry{Poisson upper}
\end{semilogyaxis}
\end{tikzpicture}
\caption{Certified bounds for $\sigma=2$, otherwise as in Fig.~\ref{fig:sigma1}. On the logarithmic axis the gap between the shuffle upper and product-pair lower curves is wider than at $\sigma=1$.}
\label{fig:sigma2}
\end{figure}

\subsection{Effect of the candidate-count refinement}
The moment/Laplace baseline bounds the two expectations in Lemma~\ref{lem:clone} by moment inequalities for $L_K$ and, in the reverse direction, also by a Laplace-transform inequality. Like the refinement, it caps each threshold at the Gaussian profile~\eqref{eq:gaussian}. The refinement never loosens that baseline at any threshold, as a comparison of the two computed profiles shows. At $\sigma=1$ it tightens the forward bound at the $411$ smallest of the $556$ thresholds and the reverse bound at the $431$ smallest; at $\sigma=2$ the counts are $247$ and $288$. Above those thresholds the two computations agree. With the same testing grid and composition routine, the five-epoch upper values decrease from $9.21$ to $7.97$ and from $0.83$ to $0.71$. Relative to the lower endpoints in Table~\ref{tab:all}, these changes remove more than $44.4\%$ and $24.3\%$ of the corresponding bracket widths. The supplementary material keeps both computations as an ablation.

At $\sigma=1$ and one epoch, the refined upper value $5.64$ falls below the deterministic value $5.77610$, which the moment/Laplace evaluation ($5.78$) did not. Since the continuous-pair epsilon exceeds $5.106$, the unrestricted optimum is at most $5.64/5.106<1.105$ times that benchmark at one epoch. The analogous two-epoch factor is $6.49/5.538<1.172$. At five epochs the certified factor is below $1.241$; for $\sigma=2$ it is below $2.101$.

\section{Lean formalization}\label{sec:lean}
\begin{table*}[t]
\centering
\caption{Paper claims, or the formalized steps of their proofs, and the corresponding Lean declarations.}
\label{tab:lean}
\footnotesize
\begin{tabular}{@{}>{\raggedright\arraybackslash}p{0.25\textwidth}>{\raggedright\arraybackslash}p{0.7\textwidth}@{}}
\toprule
Paper & Lean declarations\\
\midrule
Theorem~\ref{thm:main}, first claim & \texttt{envelope\_pair\_exists}, \texttt{envPair\_spec}, \texttt{envPair\_tradeoff}, \texttt{envPair\_dominates}\\
Shape of $f$; $T_{P,Q}\ge f$ & \texttt{envelope\_convexOn}, \texttt{envelope\_antitoneOn}, \texttt{envelope\_le\_one\_sub}, \texttt{envelope\_one}, \texttt{envelope\_le\_tradeoff}\\
Randomized tests & \texttt{integral\_sub\_le\_hockeyStick}\\
Reverse divergence (Section~\ref{sec:model}) & \texttt{hockeyStick\_swap}, \texttt{hockeyStick\_swap\_le}\\
Lemma~\ref{lem:clone}: TV transfer & \texttt{hockeyStick\_log\_le\_hockeyStick\_add\_tv}\\
Lemma~\ref{lem:clone}: modified kernel & \texttt{integral\_cloneResidual}, \texttt{isProbDensity\_cloneDensity}, \texttt{tvDist\_withDensity\_cloneDensity}, \texttt{withDensity\_cloneDensity}\\
Lemma~\ref{lem:clone}: sequential coupling & \texttt{tvDist\_partialTraj\_comp\_le\_card}\\
Lemma~\ref{lem:clone}: averaging, data processing & \texttt{hockeyStick\_sum\_smul\_le}, \texttt{hockeyStick\_comp\_le}, \texttt{hockeyStick\_map\_le}\\
Lemma~\ref{lem:clone}: $H$ in likelihood-ratio form & \texttt{hockeyStick\_log\_eq\_integral\_rnDeriv}, \texttt{hockeyStick\_swap\_log\_eq\_integral\_rnDeriv}\\
\eqref{eq:gamma}, $\eta_u\le\gamma_\sigma(p)$ & \texttt{Gaussian.integral\_clone\_eq\_cloneCost}, \texttt{Gaussian.integral\_clone\_euclidean\_le\_cloneCost}\\
\eqref{eq:gaussian} & \texttt{Gaussian.hockeyStick\_gaussianReal}, \texttt{Gaussian.hockeyStick\_gaussianReal\_symm}, \texttt{Gaussian.integral\_lr\_sub\_pos}, \texttt{Gaussian.integral\_one\_sub\_mul\_lr\_pos}\\
\eqref{eq:recurrences} & \texttt{Gaussian.callF\_succ}, \texttt{Gaussian.putJ\_succ}, \texttt{Gaussian.convexOn\_callF}, \texttt{Gaussian.convexOn\_putJ}, \texttt{callF\_eq\_callLaw}, \texttt{putJ\_eq\_putLaw}\\
Proposition~\ref{prop:counts}, given Lemma~\ref{lem:clone} & \texttt{Gaussian.counts\_bound\_forward}, \texttt{Gaussian.counts\_bound\_reverse}\\
\eqref{eq:tau} and its correction & \texttt{Gaussian.integral\_lr\_sub\_truncLR}, \texttt{Gaussian.truncLR\_le}, \texttt{call\_tail\_correction\_tau\_div}, \texttt{put\_tail\_correction}\\
\eqref{eq:spread}, Poisson spread & \texttt{spread\_masses}, \texttt{integral\_le\_spread}, \texttt{poisson\_spread}, \texttt{poisson\_spread\_convex\_order}\\
Directed convolution & \texttt{callFin\_le\_nodes}, \texttt{putFin\_le\_nodes}, \texttt{certificate\_call\_sound}, \texttt{certificate\_put\_sound}\\
Relocation of rounded mass & \texttt{expect\_le\_relocate\_top}, \texttt{expect\_le\_relocate\_bot}\\
\eqref{eq:upperpld}; lower products & \texttt{hockeyStick\_pi\_le\_upperProfile}, \texttt{lowerProfile\_le\_hockeyStick\_pi}, \texttt{isLowerSub\_prod\_iterate}\\
Integer packing & \texttt{natConv\_le\_pow\_56}, \texttt{natConv\_le\_pow\_72}, \texttt{pack128\_digit}, \texttt{pack192\_digit}\\
Appendix~\ref{app:reflection} & \texttt{Reflection.chua\_pair\_strictly\_improvable}, \texttt{Reflection.hockeyStick\_chua\_lt\_adapt}\\
Decimal ratios in the text & \texttt{decimal\_claims}\\
\bottomrule
\end{tabular}
\end{table*}

We formalized the dominating-pair construction, the counterexample and the supporting lemmas listed in Table~\ref{tab:lean} in Lean~4 with Mathlib (toolchain v4.34.0-rc2, Mathlib revision 2631d1cc). The development has seven modules and about 5500 lines, all in the namespace \texttt{ASGA}. Each of the 200 results printed by the accompanying audit file compiles without \texttt{sorry}, and \texttt{\#print axioms} reports only \texttt{propext}, \texttt{Classical.choice} and \texttt{Quot.sound}. The hockey-stick divergence is defined as a supremum over measurable sets and the testing tradeoff as an infimum over measurable tests with values in $[0,1]$. Both agree with the usual definitions for finite measures and levels $a\in[0,1]$, which is where the theorems use them, and the envelope, transfer and data-processing results hold for probability measures on arbitrary measurable spaces.

The first claim of Theorem~\ref{thm:main} is proved as stated. For every finite list of rational lines with $k_j\ge1$ and $0\le d_j^\pm\le1$, an explicit pair with rational masses has tradeoff $f$ on $[0,1]$ and dominates, in both directions, every pair of probability measures that satisfies the line inequalities. The formal construction cuts $[0,1]$ at every point where two of the lines in~\eqref{eq:envelope}, the zero line included, cross, so one segment of the envelope may be split into several atoms with the same likelihood ratio; the tradeoff is unchanged. The identity of Section~\ref{sec:model} relating the two orders, which turns forward domination at every real threshold into reverse domination, is also proved. Appendix~\ref{app:reflection} is proved for every $T\ge2$ and $\sigma>0$ on the translated transcript densities, with the forward inequality at every real $\eps$ and the reverse inequality at every $\eps\ge0$.

For Lemma~\ref{lem:clone}, the formal results are the transfer cost $(1+k)\beta_p$, the modified kernel with $\TV(g_u,\widetilde g_u)=\eta_u$, the bound $\eta_u\le\gamma_\sigma(p)$ for every $u$ in the unit ball of $\R^d$ with equality~\eqref{eq:gamma} at a unit scalar displacement, the sequential-coupling bound for $T$ history-dependent kernels (proved by a hybrid argument), data processing, and the expression of both divergences through a likelihood ratio. Proposition~\ref{prop:counts} is derived from the two conclusions of Lemma~\ref{lem:clone} and the identities $\E(X-k)_+=\E(1-kX)_+=g_\sigma(\log k)$; the certified values $\overline F_m(mk)$ and $\overline J_m(m/k)$ enter as hypotheses. For Appendix~\ref{app:certificate}, Lean checks the recurrences~\eqref{eq:recurrences} and the convexity of $F_m$ and $J_m$, the convex-order spreads, the chord induction with its boundary cases and the relocation of rounded-away mass. It also checks the closed form~\eqref{eq:tau} and, in a separate statement, the Lipschitz correction that $\tau$ pays for, together with the profile bound~\eqref{eq:upperpld}, its lower counterpart for products of lower submeasures, and carry-free integer packing. A final lemma shows that the functions $F_m$ and $J_m$ of Proposition~\ref{prop:counts} are exactly the recursions these results bound.

Several steps rest on the written proofs alone. The reduction of the shuffled mechanism to the conditional experiment $(P_K,Q_K)$, which consists of the construction in Section~\ref{sec:singleton}, the coin realization that hides the differing item among $K$ candidates and the simulation with feedback, is not formalized, and neither is the identification of the likelihood ratio of that experiment with $L_K$. Lemma~\ref{lem:clone} is therefore assembled by hand from checked parts. The following are proved only in the text: the cap of each conditional epoch by $g_\sigma$ in Proposition~\ref{prop:counts}; the adaptive composition theorem~\cite[Thm.~10]{zhu} behind the second claim of Theorem~\ref{thm:main}; the attainment of $(P_S^{\otimes E},Q_S^{\otimes E})$ and the distribution~\eqref{eq:maxcdf}; the passage from the per-bin convex-order inequality to domination of the discretized Poisson experiment with its tail atoms; the Gaussian composition behind the deterministic brackets; and, in Appendix~\ref{app:reflection}, the passage from batch sums to the translated densities. For the lower endpoints, Lean checks the data-processing step and the lower profile of products. Finally, the formal development does not model binary64 or MPFR rounding. Computed quantities enter its theorems only as inequality hypotheses, and the programs that supply them rely on the directed rounding of Appendix~\ref{app:certificate}.

\section{Limitations}
Neither endpoint of the brackets in Table~\ref{tab:all} is known to be tight. The coordinate maximum discards information, so the lower curve sits below the epsilon of the continuous product pair, and the upper accountant may be loose. The results therefore do not identify an optimal adaptive policy, and they leave open how much an adaptive analyst gains over the product pair.

The guarantee assumes fresh independent permutations, zero-out adjacency and ideal Gaussian releases. Reused permutations, replacement of one real item by another, released batch identities and a particular finite-precision noise sampler each need further analysis. The supplied profiles cover $T=1000$ and $\sigma\in\{1,2\}$, although the reduction permits other parameters and products of different epoch pairs. All bounds concern the full transcript of noisy query answers; a released final model, as a postprocessing, can have smaller privacy loss.

\section*{Open Science}
The code and certificates are provided with this submission as supplementary material. They include the rational dominating pairs, an accountant with an explicit sampler argument, every numerical input, and the programs that rebuild Table~\ref{tab:all} and both figures. A full rebuild recomputes the conditional convolutions, checks the two divergence directions and writes exact endpoint fractions; on one x86-64 core it takes under three minutes. A separate script compares every number quoted in this paper with the regenerated outputs. The Lean~4 development described in Section~\ref{sec:lean} is part of the same supplementary material. No training dataset is needed.

\section*{LLM usage considerations}
This paper was developed with AI assistance from GPT-6 Astra Pro and Claude models.

\bibliographystyle{IEEEtran}
\bibliography{references}

\appendices
\section{Certificate machinery}\label{app:certificate}
\subsection{A finite convex upper law for the likelihood ratio}
The normal input is divided into bins of width $1/50$ on $[-10,10]$. Let $X=e^{Z/\sigma-1/(2\sigma^2)}$, write $X(z)$ for its value at $Z=z$, and let $c$ be $X(10)$ rounded down to a multiple of $2^{-30}$. Define $Y=X$ on $[-10,10]$, $Y=0$ below this interval and $Y=c$ above it. Then $Y\le X$ and
\begin{align}\label{eq:tau}
 \tau=\E(X-Y)
 &=\Phi(-10-1/\sigma)+\Phi(1/\sigma-10)\nonumber\\
 &\quad-c\Phi(-10).
\end{align}

The truncation is charged explicitly. The positive part is $1$-Lipschitz, so $F_m$ computed from $Y$ needs a forward correction $m\tau$, which becomes $\tau$ after division by $m$. Since $Y\le X$ and $J_m$ is nonincreasing in each summand, the reverse put needs no correction.

For a finite normal bin $(z_0,z_1]$, the probability and first likelihood moment are
\[
 \rho=\Phi(z_1)-\Phi(z_0),\qquad
 \mu=\Phi(z_1-1/\sigma)-\Phi(z_0-1/\sigma).
\]
The support endpoints $l\le X(z_0)$ and $v\ge X(z_1)$ are taken on the grid $2^{-30}\mathbb Z$, rounding outward.

The bin is replaced by the endpoint masses
\begin{equation}\label{eq:spread}
 w_l=\frac{v\rho-\mu}{v-l},\qquad
 w_v=\frac{\mu-l\rho}{v-l}.
\end{equation}
These masses are nonnegative, sum to $\rho$ and preserve the bin's first moment. By the chord inequality, any convex function of the bin's value has expectation at most its expectation under the two endpoint masses, so the spread law dominates $Y$ in convex order. The two tail atoms of $Y$ are kept unchanged.

Endpoint probabilities are then rounded downward to multiples of $2^{-72}$, which leaves a small unassigned mass. The forward upper law moves all of it to the largest support point, which can only increase $F_m(t)$, since $F_m$ is nondecreasing in each summand. The reverse upper law moves it to zero, which can only increase $J_m(t)$, since $J_m$ is nonincreasing in each summand. Both laws are normalized exactly.

\subsection{Directed conditional convolution}
Threshold nodes are dyadic. In units of $2^{-30}$ the grid starts at $0$ and $1$ and advances by the larger of $1$ and the current node divided by $400$ and rounded down, until it reaches the required maximum. Forward nodes extend through $M$ times the largest support point and reverse nodes through $M$. Node placement affects accuracy only, since the chord argument below is valid on any increasing grid.

Each function in~\eqref{eq:recurrences} is computed from its predecessor. Suppose the exact $F_{m-1}$ is convex and bounded above at the nodes by stored values. At a point $t-y$ inside a grid interval, convexity bounds $F_{m-1}(t-y)$ by the chord through the exact endpoint values, and hence by the chord through the stored values. Averaging these chords over the atoms $y$ of the upper law bounds $F_m(t)$, and the same argument applies to $J_m$.

The stored values bound $F_m$ and $J_m$ computed with the upper laws as summand laws. These functions dominate those for $Y$, because the upper laws dominate $Y$ in increasing and decreasing convex order respectively, and $F_m$ and $J_m$ are convex and monotone in each summand.

Two boundary cases have closed forms. For $t-y<0$ the forward continuation is at most $(m-1)\bar y-(t-y)$, where $\bar y$ is an upward-rounded mean of the forward upper law, because $S_{m-1}\ge0$; the reverse continuation is zero. Beyond the largest possible forward sum the forward continuation is zero. Together with the exact base functions $F_0$ and $J_0$ these cases complete the induction.

All transition arithmetic is nonnegative and rounded toward $+\infty$ in IEEE binary64. Integer support differences are exact, and each probability is padded upward by one representable step when it is converted to binary64. Compilation disables contraction and fast-math transformations, and the executable checks the active rounding mode when it starts and when it exits.

The value $F_m(mk)$ is read at an argument rounded downward, which is safe because $F_m$ is nonincreasing in $t$; $J_m(m/k)$ is read at an argument rounded upward because $J_m$ is nondecreasing in $t$. The tail correction~\eqref{eq:tau} is added to the forward value, and each conditional profile is capped by~\eqref{eq:gaussian}.

Floating-point search only proposes the clone parameter, from the grids $\{140,\ldots,380\}/40$ for $\sigma=1$ and $\{65,\ldots,240\}/40$ for $\sigma=2$. The accepted expression~\eqref{eq:finite-count} is then re-evaluated with directed intervals over all $m\le M$, using an upper enclosure for the residual binomial probability.

As a check on discretization error, the whole upper computation was repeated with $25$ normal bins per unit and grid denominator $200$. It certifies the same epsilon values at $\sigma=1$, and values larger by at most $0.02$ at $\sigma=2$.

\subsection{Rational testing envelopes and upper products}
GNU MPFR~\cite{mpfr} evaluates elementary functions and the complementary error function with separate downward and upward rounding at $256$-bit precision. The interval wrapper exports exact dyadic endpoints, and every completed one-epoch bound is rounded upward to a multiple of $2^{-100}$ and capped at one.

Sorting the rational lines of~\eqref{eq:envelope} by slope gives their upper envelope, whose slopes and intersections yield the atoms~\eqref{eq:atoms}. The verifier checks that the envelope lies above every input line, that the masses are nonnegative and exactly normalized, and that every finite atom has the stated likelihood ratio.

For composition, the finite masses are rounded downward to multiples of $2^{-56}$. The resulting arrays are lower submeasures at known likelihood ratios, and they are never renormalized.

Integer polynomial multiplication computes each convolution exactly, after which every coefficient is divided by $2^{56}$ and rounded down. Packing the coefficients in base $2^{128}$ is safe because each product coefficient is at most $2^{112}$. After unpacking, the program checks that the coefficient sum equals the product of the input sums.

For a retained real-law mass $\underline w_j$ at privacy loss $j\log r$, a certified upper profile is
\begin{equation}\label{eq:upperpld}
 1-\sum_j\underline w_j\min\{1,e^\eps r^{-j}\}.
\end{equation}
Every missing mass, including the singular atoms and all rounding deficits, is thereby paid at infinite loss. The factors subtracted in~\eqref{eq:upperpld} are rounded downward.

The second direction uses the null-law array with reversed loss indices in the same formula. The deterministic Gaussian brackets come from an independent rational bisection on~\eqref{eq:gaussian} with $\sigma$ replaced by $\sigma/\sqrt E$.

\subsection{Lower products from the maximum statistic}
For each bin of~\eqref{eq:maxcdf}, let $p_i$ and $q_i$ be its exact probabilities. Directed intervals enclose both probabilities and $\log(p_i/q_i)$. Under the real law, masses are rounded downward to multiples of $2^{-72}$ and privacy losses downward to multiples of $\log(2001/2000)$; under the null law the same is done for the reversed loss $\log(q_i/p_i)$.

Because $\ell\mapsto(1-e^{\eps-\ell})_+$ is nondecreasing, products of these rounded-down losses with the missing mass discarded give lower bounds on the divergences. Rounding losses upward and paying missing mass at infinite loss gives upper bounds for the binned experiment itself, which supply the inverse intervals in Section~\ref{sec:evaluation}. Both products use exact integer convolution, packed in base $2^{192}$ for the $72$-bit masses.

\subsection{A continuous Poisson upper comparison}
Under $B$ in~\eqref{eq:poisson}, the likelihood ratio is
\[
 L(z)=1-q+qe^{z/\sigma-1/(2\sigma^2)},\qquad z\sim\N(0,1).
\]
The likelihood axis is cut at the points $r^j$ with $r=10001/10000$.

For a likelihood bin $[\lambda,\lambda']=[r^j,r^{j+1}]$ with $B$-mass $q_i$ and $A$-mass $p_i$, the endpoint $B$-masses are
\[
 s_\lambda=\frac{\lambda'q_i-p_i}{\lambda'-\lambda},\qquad
 s_{\lambda'}=\frac{p_i-\lambda q_i}{\lambda'-\lambda},
\]
and the endpoint $A$-masses are $\lambda s_\lambda$ and $\lambda's_{\lambda'}$. For any convex function of $L$ the endpoint chord bounds its conditional expectation, so the resulting binary experiment dominates $(A,B)$ and adaptive composition applies to it.

The bin probabilities are explicit Gaussian CDF differences. For $u>1-q$ the likelihood ratio reaches $u$ at
\[
 z(u)=\sigma\left[\log\frac{u-1+q}{q}+\frac1{2\sigma^2}\right],
\]
and the two cumulative probabilities are $\Phi(z(u))$ under $B$ and $(1-q)\Phi(z(u))+q\Phi(z(u)-1/\sigma)$ under $A$.

The lowest likelihood node is at most $1-q$, the infimum of $L$. The final node lies beyond $L(10)$, and the remaining tail is split into hypothesis-specific singular atoms. Revealing that tail perfectly can only increase the divergences, so domination is preserved.

Retained endpoint masses are rounded downward to multiples of $2^{-72}$ and composed for $ET$ steps with exact integer multiplication. Cropping after intermediate convolutions is allowed because the whole deficit is charged in~\eqref{eq:upperpld}, and separate real and reverse-null arrays cover both directions. The result is an upper bound for the continuous mechanism with no unquantified quadrature error.

\section{The proposed pair is strictly improvable}\label{app:reflection}
The differing item can change its query value in response to a prefix that favors the real-input hypothesis. Such a response can remove cancellation in the final likelihood-ratio test. The following argument also explains why an attainable pair need not be a safe adaptive accountant.

Fix $T\ge2$, $\sigma>0$, $k=e^\eps$, and common query values $-1$. The differing value is $+1$ in the first $T-1$ rounds. With $z_t=(Y_t+b)/\sigma$, set
\[
 U=\sum_{t<T}e^{2z_t/\sigma-2/\sigma^2},\qquad
 V=\sum_{t<T}e^{z_t/\sigma-1/(2\sigma^2)}.
\]
Change the final differing value to $-1$ precisely when $r=U-kV>0$, leaving the null value zero.

Write $f_u$ for the $\N(u,\sigma^2)$ density, with $f_\pm=f_{\pm1}$. In the final round use the coordinate $Y_T+b-1$. Its mean is $1$ if the real differing item is in that batch, $0$ if its null replacement is, and $-1$ if the batch holds only common items. At a fixed prefix, both hypotheses' densities share the positive factor $T^{-1}\prod_{t<T}f_0(Y_t+b)$, which we drop; $U$ and $V$ then collect the likelihood ratios of the prefix rounds. The forward signed densities before and after changing the query are $rf_-+f_+-kf_0$ and $(r+1)f_--kf_0$. Convexity gives
\begin{align*}
 \int(rf_-+f_+-kf_0)_+
 &\le\frac r{r+1}\int((r+1)f_--kf_0)_+\\
 &\quad+\frac1{r+1}\int((r+1)f_+-kf_0)_+\\
 &=\int((r+1)f_--kf_0)_+.
\end{align*}
The endpoint integrals agree by reflection. Their signed densities have opposite half-line positive regions and disagree in sign on a positive-measure set. Both mixture weights are positive, so the inequality is strict.

For the reverse direction put $s=V-kU$. On $r>0$ and $k\ge1$, $s<0$. The baseline density difference is $f_0-kf_+-(-s)f_-$; after switching it is $f_0-(k-s)f_-$. The baseline difference is the mixture of $f_0-(k-s)f_+$ and $f_0-(k-s)f_-$ with weights $k/(k-s)$ and $(-s)/(k-s)$, both positive. The same strict convexity argument and a reflection show that switching strictly increases the integral of the positive part.

The switching region has positive probability: $U/V\to\infty$ as $z_1\to\infty$ with the other prefix coordinates fixed, so $\{U>kV\}$ is a nonempty open subset of $\R^{T-1}$ on which the prefix has positive density under both hypotheses. Outside it the mechanisms coincide. Integrating shows that the forward divergence increases strictly at every real $\eps$, since that argument needs only $k>0$, and the reverse divergence at every $\eps\ge0$. The switching policy depends on the tested $\eps$, so the argument does not produce one policy that exceeds the pair at every $\eps$.
\end{document}
